%! AP Statistics — Unit 2, Lesson 2.4 — Group Activity
\documentclass[10pt]{article}
\usepackage{apstats-article}
\usepackage{apstats-boxes}

%=============================================================================
\begin{document}

\pageheader{Unit 2, Lesson 2.4}{Group Activity}
\namepartnerperiod

\vspace{0.18in}

% ── SCENARIO ─────────────────────────────────────────────────────────────────
\begin{scenariobox}[Context]{sky}
\small
Your group has been given a real-world dataset. Your task is to construct a scatterplot,
describe the relationship using the DOFS framework, and draw a conclusion in context.
Each tier uses a different dataset --- compare results with other groups at the end.
\end{scenariobox}

\vspace{0.12in}

% ── TIER R ────────────────────────────────────────────────────────────────────
\begin{tcolorbox}[colback=white, colframe=black!40, boxrule=0.8pt, arc=2mm,
    left=4mm, right=4mm, top=3mm, bottom=3mm,
    title={\textbf{Tier R --- Remediate}}, fonttitle=\bfseries, halign title=center]
\small

\textbf{Dataset: Study time vs.\ quiz score (12 students)}

\vspace{6pt}
\renewcommand{\arraystretch}{1.4}
\begin{tabularx}{\linewidth}{@{} l *{12}{C{0.72cm}} @{}}
\rowcolor{sky}
\textbf{Study (hr)} & 0 & 0.5 & 1 & 1 & 1.5 & 2 & 2 & 2.5 & 3 & 3 & 3.5 & 4 \\
\textbf{Score}      & 55 & 60 & 62 & 65 & 68 & 72 & 75 & 78 & 80 & 84 & 86 & 90 \\
\end{tabularx}

\vspace{10pt}

\begin{enumerate}[leftmargin=1.6em, itemsep=0.18in]

  \item Plot each point on the grid below. Label both axes (variable name and units).

        \vspace{6pt}
        \begin{center}
        \begin{tikzpicture}
          \draw[very thin, gray!35] (0,0) grid (8,6);
          \draw[->, thick, navy] (0,0) -- (8.4,0);
          \draw[->, thick, navy] (0,0) -- (0,6.5);
          \node[below=8pt] at (4,0) {\small \blank{5cm}};
          \node[left=8pt, rotate=90] at (0,3) {\small \blank{5cm}};
        \end{tikzpicture}
        \end{center}

  \item Write a complete DOFS description in context.\\[8pt]
        \writeline\\[2pt]\writeline\\[2pt]\writeline

  \item Does any point appear to be an outlier? If so, describe it in context and speculate
        on a possible reason.\\[8pt]
        \writeline\\[2pt]\writeline

\end{enumerate}
\end{tcolorbox}

\vspace{0.12in}

% ── TIER A ────────────────────────────────────────────────────────────────────
\begin{tcolorbox}[colback=white, colframe=black!40, boxrule=0.8pt, arc=2mm,
    left=4mm, right=4mm, top=3mm, bottom=3mm,
    title={\textbf{Tier A --- Approaching Proficiency}}, fonttitle=\bfseries, halign title=center]
\small

\textbf{Dataset: Age of car (years) vs.\ resale price (\$1000s) --- 12 cars}

\vspace{6pt}
\renewcommand{\arraystretch}{1.4}
\begin{tabularx}{\linewidth}{@{} l *{12}{C{0.72cm}} @{}}
\rowcolor{sky}
\textbf{Age (yr)} & 1 & 2 & 3 & 3 & 4 & 5 & 5 & 6 & 7 & 8 & 10 & 12 \\
\textbf{Price}    & 28 & 25 & 22 & 20 & 18 & 16 & 14 & 13 & 11 & 8 & 6 & 3 \\
\end{tabularx}

\vspace{10pt}

\begin{enumerate}[leftmargin=1.6em, itemsep=0.18in]

  \item Identify the explanatory and response variables. Plot each point on the grid below.
        Label both axes completely.

        Explanatory: \blank{4cm} \quad Response: \blank{4cm}

        \vspace{6pt}
        \begin{center}
        \begin{tikzpicture}
          \draw[very thin, gray!35] (0,0) grid (8,6);
          \draw[->, thick, navy] (0,0) -- (8.4,0);
          \draw[->, thick, navy] (0,0) -- (0,6.5);
          \node[below=8pt] at (4,0) {\small \blank{5cm}};
          \node[left=8pt, rotate=90] at (0,3) {\small \blank{5cm}};
        \end{tikzpicture}
        \end{center}

  \item Write a complete DOFS description in context.\\[8pt]
        \writeline\\[2pt]\writeline\\[2pt]\writeline

  \item A 50-year-old classic car has a resale price of \$45{,}000. Plot this point. Does it
        follow the overall pattern? Explain what kind of point this is and why it is unusual.\\[8pt]
        \writeline\\[2pt]\writeline\\[2pt]\writeline

\end{enumerate}
\end{tcolorbox}

\vspace{0.12in}

% ── TIER E ────────────────────────────────────────────────────────────────────
\begin{tcolorbox}[colback=white, colframe=black!40, boxrule=0.8pt, arc=2mm,
    left=4mm, right=4mm, top=3mm, bottom=3mm,
    title={\textbf{Tier E --- Extension}}, fonttitle=\bfseries, halign title=center]
\small

\textbf{Dataset: Distance from city center (miles) vs.\ average monthly rent (\$) --- 10 neighborhoods}

\vspace{6pt}
\renewcommand{\arraystretch}{1.4}
\begin{tabularx}{\linewidth}{@{} l *{10}{C{0.82cm}} @{}}
\rowcolor{sky}
\textbf{Distance} & 0.5 & 1 & 2 & 3 & 4 & 5 & 7 & 9 & 10 & 12 \\
\textbf{Rent}     & 3200 & 2900 & 2500 & 2200 & 2000 & 1800 & 1600 & 1400 & 1300 & 800 \\
\end{tabularx}

\vspace{10pt}

\begin{enumerate}[leftmargin=1.6em, itemsep=0.18in]

  \item Construct the scatterplot on the grid below, choosing appropriate axis scales.

        \vspace{6pt}
        \begin{center}
        \begin{tikzpicture}
          \draw[very thin, gray!35] (0,0) grid (8,6);
          \draw[->, thick, navy] (0,0) -- (8.4,0);
          \draw[->, thick, navy] (0,0) -- (0,6.5);
          \node[below=8pt] at (4,0) {\small \blank{5cm}};
          \node[left=8pt, rotate=90] at (0,3) {\small \blank{5cm}};
        \end{tikzpicture}
        \end{center}

  \item Write a complete DOFS description in context.\\[8pt]
        \writeline\\[2pt]\writeline\\[2pt]\writeline

  \item The pattern looks roughly linear, but could also be slightly curved. Explain how
        you would decide which description is more accurate, and why the form matters for
        what we do next in the course.\\[8pt]
        \writeline\\[2pt]\writeline\\[2pt]\writeline

  \item A neighborhood 3 miles from the center has a monthly rent of \$3{,}500. Add this
        point to your scatterplot. Describe it in context and speculate on why rent might
        be so high despite being far from the center.\\[8pt]
        \writeline\\[2pt]\writeline

\end{enumerate}
\end{tcolorbox}

\end{document}
